Can wider search worsen deployment when its selection score omits a cost? Selection overfitting and the effect of trading costs are established concerns. We use standard finite-pool rank probabilities to separate changes in winner exposure from the measured gross-to-net outcome gap. In a trading simulation, four pools of 500--1,500 ridge-regression strategies are ranked by gross Sharpe, ignoring execution costs. Wider search selects higher turnover, and the capacity component of the gap has a positive fitted slope on the primary search grid in all 16 pool-by-capital cells. From one to 32 trials, mean net Sharpe falls in 11 of 16 cells and rises in all four pool-A cells. The capacity effect is not statistically established: its joint selection--evaluation bootstrap interval for pool B at USD 500 million is $[-0.022,0.121]$. Mean candidate gross Sharpe is negative in every pool, and all tested rules except an infeasible rule that selects on evaluation outcomes have negative mean net Sharpe on the comparison grid. Selecting on modeled net Sharpe improves that mean by $0.430$, with paired interval $[0.192,0.723]$, but leaves it at $-0.071$. The calculation makes search-induced exposure to an omitted term explicit. It also separates that measurement problem from selection noise: statistical correction cannot recover a cost the score does not measure.
